\chapter{Ruang Vektor}\label{ch:b1:vspaces}

Aljabar linear bermula di sini: aksioma \omterm{def:b1:vspaces:def}{ruang vektor} mengucilkan apa yang
dimiliki bersama oleh $\R^2$, $\R^3$, ruang \omterm{def:b1:poly:def}{polinomial} dan ruang fungsi
--- yaitu bahwa kita dapat menjumlahkan, dan menskalakan. Dua bab membangun teorinya
(dengan \cref{ch:b1:findim} menambahkan dimensi); dan bahasa yang disiapkannya ---
\omterm{def:b1:vspaces:span}{rentang}, \omterm{def:b1:vspaces:free}{keluarga bebas}, \omterm{def:b1:vspaces:free}{basis}, \omterm{def:b1:vspaces:sum}{jumlah langsung} --- merupakan santapan harian setiap
bab sesudahnya. Di sepanjang bab ini, $K$ menyatakan $\R$ atau $\C$
(yaitu \emph{skalarnya}).

\section{Definisi dan contoh}

\begin{definition}[Ruang vektor]\label{def:b1:vspaces:def}
Sebuah \emph{ruang vektor atas $K$}\index{ruang vektor} adalah \omterm{def:b1:logic:sets}{himpunan} $E$ dengan sebuah
penjumlahan yang membuat $(E, +)$ menjadi \omterm{def:b1:structures:group}{grup abelian} (dengan nolnya ditulis $0_E$ atau
$0$), beserta perkalian skalar $K \times E \to E$ sedemikian sehingga, untuk
setiap $\lambda, \mu \in K$ dan $x, y \in E$:
\[
\lambda(x + y) = \lambda x + \lambda y,\quad
(\lambda + \mu) x = \lambda x + \mu x,\quad
\lambda(\mu x) = (\lambda\mu) x,\quad
1\,x = x .
\]
Akibatnya: $0\,x = 0_E$, $\lambda\,0_E = 0_E$, $(-1)x = -x$, dan
$\lambda x = 0_E \implies \lambda = 0$ atau $x = 0_E$ (dengan mengalikannya dengan
$\lambda^{-1}$).
\end{definition}

\begin{proof}[Bukti akibatnya]
Untuk $0\,x = 0_E$: dari $(0 + 0)x = 0x + 0x$ dan $(0+0)x = 0x$,
coretlah $0x$ pada \omterm{def:b1:structures:group}{grup} $(E, +)$. Untuk $\lambda\,0_E$: dengan muslihat
yang sama pada $\lambda(0_E + 0_E)$. Untuk $(-1)x$: tambahkanlah $x$,
\[
x + (-1)x = 1\,x + (-1)x = \bigl(1 + (-1)\bigr)x = 0\,x = 0_E ,
\]
sehingga $(-1)x$ merupakan invers penjumlahan $x$. Akhirnya jika $\lambda x
= 0_E$ dengan $\lambda \neq 0$: kalikanlah dengan $\lambda^{-1}$ (karena
skalarnya membentuk \omterm{def:b1:structures:field}{lapangan}) lalu pakailah kedua aksioma
$\lambda^{-1}(\lambda x) = (\lambda^{-1}\lambda)x = 1x = x$
beserta $\lambda^{-1}0_E = 0_E$: sehingga $x = 0_E$. Sekecil apa pun
keempat aturan ini, ia dipakai diam-diam pada setiap halaman
berikutnya --- dan yang terakhir persis merupakan tempat \omterm{def:b1:structures:field}{lapangan} diperlukan:
karena di atas skalar $\Z$, ``ruang'' $\Z/2\Z$ akan melanggarnya
dengan $2\,x = 0$.
\end{proof}

\begin{example}\label{ex:b1:vspaces:examples}
$K^n$ (dengan operasi \omterm{prop:b1:vspaces:coordinates}{koordinat} demi \omterm{prop:b1:vspaces:coordinates}{koordinat}); \omterm{def:b1:poly:def}{polinomial} $K[X]$; fungsi
$\mathcal{F}(A, K)$ dari sebarang \omterm{def:b1:logic:sets}{himpunan} $A$ ke $K$ (dengan operasi
titik demi titik) --- yang memuat fungsi yang \omterm{def:b1:continuity:continuous}{kontinu}, barisan
$\mathcal{F}(\N, \R)$, dan seterusnya; serta $\C$ sebagai \omterm{def:b1:vspaces:def}{ruang vektor} atas $\R$. Pada setiap
kasusnya aksiomanya diwarisi dari aksioma $K$.
\end{example}

\begin{definition}[Subruang]\label{def:b1:vspaces:subspace}
\omterm{def:b1:logic:sets}{Himpunan} $F \subseteq E$ disebut \emph{subruang}\index{subruang} bila $0_E \in
F$ dan $F$ stabil terhadap penjumlahan dan perkalian skalar ---
setara dengan itu:
\[
F \neq \emptyset
\qquad\text{dan}\qquad
\forall x, y \in F,\ \forall \lambda \in K,\quad
x + \lambda y \in F .
\]
Sebuah subruang sendirinya merupakan \omterm{def:b1:vspaces:def}{ruang vektor}. Adapun sebarang irisan subruang
merupakan subruang; sedangkan gabungannya hampir tak pernah demikian (dengan bukti yang sama seperti
\cref{exo:b1:structures:6}).
\end{definition}

\begin{example}\label{ex:b1:vspaces:subspaces}
Di dalam $\mathcal{F}(\R, \R)$: fungsi yang \omterm{def:b1:continuity:continuous}{kontinu}, fungsi yang
\omterm{def:b1:derivative:def}{dapat diturunkan}, \omterm{def:b1:poly:def}{polinomial} berderajat $\leq n$ (yang ditulis
$K_n[X]$ di dalam $K[X]$), dan penyelesaian persamaan diferensial linear
yang homogen (\cref{thm:b1:diffeq:homogeneous2} mengatakan persis
itu). Adapun yang bukan contoh: $\{f : f(0) = 1\}$ (karena tanpa nol); dan berderajat tepat
$n$ (karena tak stabil terhadap penjumlahan).
\end{example}

\begin{example}[Subruang atau bukan: empat vonis, beserta argumennya]
\label{ex:b1:vspaces:subspaceornot}
Di dalam ruang barisan real:
\begin{itemize}
  \item $\{u : u \text{ terbatas}\}$ \emph{memang} \omterm{def:b1:vspaces:subspace}{subruang}: karena $0$
        terbatas, dan jika $\abs{u_n} \leq M$, $\abs{v_n} \leq M'$,
        maka $\abs{u_n + \lambda v_n} \leq M + \abs\lambda M'$.
  \item $\{u : u_n \to 1\}$ \emph{bukan}: karena barisan nolnya
        hilang (dan jumlah dua anggotanya menuju $2$).
  \item $\{u : u \text{ monoton}\}$ \emph{bukan}: karena $u_n = n$
        dan $v_n = -n + (-1)^n$ monoton, sedangkan jumlahnya
        $(-1)^n$ tidak; jadi kestabilan terhadap penjumlahanlah aksioma
        yang gagal, meskipun \omterm{def:b1:logic:sets}{himpunannya} memuat $0$ dan semua
        kelipatan skalar anggotanya.
  \item $\{u : u_{n+1} = u_n^2\}$ \emph{bukan}: karena ia memuat $0$
        tetapi $2u$ kabur begitu $u$ menjadi anggota yang taknol
        (sebab $2u_{n+1} \neq (2u_n)^2$ pada umumnya) --- jadi pengkuadratannya
        yang menjadi ketaklinearannya.
\end{itemize}
Adapun urutan kerjanya selalu sama: ujilah $0$ lebih dulu (karena paling murah),
lalu kestabilannya --- dan untuk membantahnya, satu pasang contoh penyangkal
yang eksplisit mengalahkan keraguan sebanyak apa pun.
\end{example}

\section{Rentang, jumlah, jumlah langsung}

\begin{definition}[Kombinasi linear, rentang]\label{def:b1:vspaces:span}
Sebuah \emph{kombinasi linear} keluarga $(x_1, \dots, x_p)$ berisi
vektor $E$ adalah sebarang $\lambda_1 x_1 + \dots + \lambda_p x_p$
(dengan $\lambda_i \in K$). Adapun \omterm{def:b1:logic:sets}{himpunan} semuanya merupakan
\emph{rentang}\index{rentang}nya $\operatorname{Vect}(x_1, \dots, x_p)$: yang
merupakan \omterm{def:b1:vspaces:subspace}{subruang}, yaitu yang terkecil yang memuat keluarganya.
\end{definition}

\begin{proof}[Bukti kedua penegasannya]
Kestabilannya: jumlah dua kombinasi linear $\sum\lambda_i x_i +
\sum\mu_i x_i = \sum(\lambda_i + \mu_i)x_i$ kembali menjadi satu, dan demikian
pula kelipatan skalarnya $\mu\sum\lambda_i x_i = \sum(\mu\lambda_i)
x_i$; sedangkan kombinasi nolnya menunjukkan bahwa $0$ termasuk: jadi \omterm{def:b1:vspaces:span}{rentangnya} sebuah
\omterm{def:b1:vspaces:subspace}{subruang}. Adapun keminimalannya: misalkan $H$ sebarang \omterm{def:b1:vspaces:subspace}{subruang} yang memuat $x_1,
\dots, x_p$. Maka menurut kestabilan terhadap perkalian skalar setiap
$\lambda_i x_i \in H$, dan menurut kestabilan terhadap penjumlahan jumlahnya
terletak di $H$: sehingga setiap kombinasi linearnya termasuk $H$, yakni
$\operatorname{Vect}(x_1, \dots, x_p) \subseteq H$. Jadi \omterm{def:b1:vspaces:span}{rentangnya}
termuat di setiap \omterm{def:b1:vspaces:subspace}{subruang} yang memuat keluarganya: sehingga ia yang
terkecil.
\end{proof}

\begin{definition}[Jumlah, jumlah langsung]\label{def:b1:vspaces:sum}
Untuk \omterm{def:b1:vspaces:subspace}{subruang} $F, G$ pada $E$:
\[
F + G = \{\,u + v : u \in F,\ v \in G\,\}
\]
merupakan \omterm{def:b1:vspaces:subspace}{subruang} (yaitu yang terkecil yang memuat $F \cup G$). Jumlahnya disebut
\emph{langsung}\index{jumlah langsung}, yang ditulis $F \oplus G$, bila setiap
unsur $F + G$ terurai \emph{secara tunggal} sebagai $u + v$;
setara dengan itu (lihatlah di bawah) bila $F \cap G = \{0\}$. Adapun ketika $E = F \oplus
G$, kedua \omterm{def:b1:vspaces:subspace}{subruangnya} disebut \emph{saling melengkapi}\index{subruang yang saling
melengkapi} di dalam $E$.
\end{definition}

\begin{example}[Sebuah jumlah dua garis]
\label{ex:b1:vspaces:sumoflines}
Di dalam $\R^3$, misalkan $F = \operatorname{Vect}\bigl((1,0,1)\bigr)$ dan
$G = \operatorname{Vect}\bigl((0,1,1)\bigr)$. Maka jumlahnya adalah
\[
F + G = \{\,a(1,0,1) + b(0,1,1)\,\}
= \{(a,\ b,\ a + b)\} = \{(x, y, z) : z = x + y\},
\]
yaitu bidang lewat titik asal yang memuat kedua garisnya. Ia
tegas lebih besar daripada gabungannya $F \cup G$ (yang sekadar silang kedua
garisnya): karena vektor $(1, 1, 2) = (1,0,1) + (0,1,1)$ terletak di
jumlahnya tetapi tak pada garis mana pun. Dan $F \cap G = \{0\}$ (karena sebuah vektor
bersama menuntut $a(1,0,1) = b(0,1,1)$, yang dua \omterm{prop:b1:vspaces:coordinates}{koordinat}
pertamanya memaksa $a = b = 0$): sehingga jumlahnya langsung, dan $F \oplus
G$ persis merupakan bidang itu.
\end{example}

\begin{proposition}\label{prop:b1:vspaces:directsum}
Jumlah $F + G$ bersifat langsung jika dan hanya jika $F \cap G = \{0\}$.
\end{proposition}

\begin{proof}
Jika suatu $w \neq 0$ terletak di $F \cap G$: maka $w = w + 0 = 0 + w$ merupakan dua
penguraian $w$. Sebaliknya, jika $u + v = u' + v'$ dengan $u, u'
\in F$, $v, v' \in G$, maka $u - u' = v' - v$ termasuk $F \cap G =
\{0\}$: sehingga penguraiannya tunggal.
\end{proof}

\begin{method}[Membuktikan $E = F \oplus G$]
\label{met:b1:vspaces:directsum}
Ada dua hal yang harus diperiksa, masing-masing dengan langkah pembuka bakunya.
\begin{enumerate}
  \item \emph{Irisan yang sepele.} Ambillah $x \in F \cap G$,
        tuliskanlah kedua syarat keanggotaannya, lalu apitlah $x =
        0$. (Jangan pernah berargumen lewat gambar: bandingkanlah jebakan di bawah.)
  \item \emph{Jumlahnya segalanya.} Ambillah sebarang $x \in E$
        lalu \emph{hasilkanlah} penguraiannya $x = f + g$ ---
        entah dengan menebak $f$ dari sasarannya (karena $f$ harus memenuhi
        sifat pendefinisi $F$, yang biasanya mendiktekan
        rumusnya) entah dengan memecahkan sistem linear yang menyatakan $x$
        terhadap pembangun $F$ dan $G$.
\end{enumerate}
Ketika rumus penguraiannya ditebak, ketunggalannya
otomatis dari langkah 1; sedangkan ketika hanya keberadaannya yang tak jelas, langkah 2-lah
tempat kerjanya tinggal. Adapun kedua contoh di bawah menjalankan metodenya: untuk
fungsi genap/ganjil rumus $f$-nya dipaksa oleh penilaian kesamaan yang
diinginkan di $x$ dan $-x$; sedangkan untuk \omterm{def:b1:poly:def}{polinomial} yang lenyap
di sebuah titik, oleh penilaian di $a$.
\end{method}

\begin{example}\label{ex:b1:vspaces:evenodd}
Di dalam $\mathcal{F}(\R, \R)$, fungsi genap $\mathcal{P}$ dan
fungsi ganjil $\mathcal{I}$ \omterm{def:b1:vspaces:sum}{saling melengkapi}: karena sebarang $f$ tertulis
\[
f(x) = \underbrace{\frac{f(x) + f(-x)}{2}}_{\text{genap}}
+ \underbrace{\frac{f(x) - f(-x)}{2}}_{\text{ganjil}},
\]
dan sebuah fungsi yang sekaligus genap dan ganjil bernilai nol. (Bila diterapkan pada $\exp$, ini
merupakan pasangan $(\cosh, \sinh)$ pada \cref{ch:b1:functions}.)
\end{example}

\begin{example}[{Sepasang pelengkap di dalam $K_n[X]$}]
\label{ex:b1:vspaces:hyperplane}
Tetapkanlah $a \in K$ lalu tetapkan $F = \{P \in K_n[X] : P(a) = 0\}$, $G =
\operatorname{Vect}(1)$ (yaitu konstantanya). Maka $K_n[X] = F \oplus
G$. Memang $F \cap G$ terdiri atas konstanta yang lenyap di $a$,
yakni $\{0\}$; dan setiap $P$ terurai sebagai
\[
P = \underbrace{\bigl(P - P(a)\bigr)}_{\in F}
+ \underbrace{P(a)}_{\in G} .
\]
Penguraian ini pantas dihafalkan: karena mengurangkan nilainya di sebuah
titik merupakan cara baku memproyeksikan ke ``fungsi yang lenyap di
$a$''. Perhatikanlah bahwa $F$ \omterm{def:b1:vspaces:subspace}{subruang} yang besar sedangkan $G$ yang kecil; jadi
sepasang \omterm{def:b1:vspaces:sum}{pelengkap} tak harus seimbang dalam arti apa pun.
\end{example}

\begin{example}[Sebuah subruang pelengkap tak pernah tunggal]
\label{ex:b1:vspaces:manysupplements}
Di dalam $\R^2$, misalkan $F = \operatorname{Vect}\bigl((1,0)\bigr)$ (yaitu
sumbu $x$). Maka baik $G = \operatorname{Vect}\bigl((0,1)\bigr)$ maupun
$G' = \operatorname{Vect}\bigl((1,1)\bigr)$ melengkapi
$F$: karena masing-masingnya bertemu $F$ hanya di $0$, dan setiap pasangannya berjumlah $\R^2$.
Adapun penguraian vektor yang sama itu berbeda:
\[
(2,\ 1.5) = \underbrace{(2, 0)}_{\in F} +
\underbrace{(0, 1.5)}_{\in G}
= \underbrace{(0.5,\ 0)}_{\in F} +
\underbrace{(1.5,\ 1.5)}_{\in G'} .
\]
Sesungguhnya \emph{setiap} garis selain $F$ sendiri merupakan
\omterm{def:b1:vspaces:sum}{pelengkap} $F$ di dalam $\R^2$: jadi \omterm{def:b1:vspaces:sum}{pelengkapnya} berlimpah,
dan berbicara tentang ``sang'' \omterm{def:b1:vspaces:sum}{pelengkap} tak bermakna sampai sebuah
struktur tambahan (yaitu hasil kali dalam, \cref{ch:b1:euclid}) memilih
salah satunya.
\end{example}

\begin{omfigure}
\begin{tikzpicture}[scale=1.5]
  \draw[->, thin] (-0.4,0) -- (2.9,0) node[below] {$x$};
  \draw[->, thin] (0,-0.4) -- (0,1.9) node[left] {$y$};
  \draw[omDef, very thick] (-0.3,0) -- (2.7,0)
    node[below, pos=0.95, font=\small] {$F$};
  \draw[omThm, very thick] (0,-0.3) -- (0,1.7)
    node[left, pos=0.95, font=\small] {$G$};
  \draw[omProp, very thick] (-0.25,-0.25) -- (1.75,1.75)
    node[right, pos=0.97, font=\small] {$G'$};
  \fill (2,1.5) circle (0.045) node[above right, font=\small]
    {$(2, 1.5)$};
  \draw[omThm, dashed] (2,1.5) -- (2,0);
  \draw[omDef, dashed] (2,1.5) -- (0,1.5);
  \draw[omProp, dashed] (2,1.5) -- (0.5,0);
  \fill (2,0) circle (0.035);
  \fill (0.5,0) circle (0.035);
\end{tikzpicture}

{\small Dua penguraian titik yang sama pada $\R^2$ sepanjang $F$
(yaitu sumbu $x$): dengan \omterm{def:b1:vspaces:sum}{pelengkap} $G$ (lewat jatuhan tegak) dan dengan
\omterm{def:b1:vspaces:sum}{pelengkap} $G'$ (lewat jatuhan miring). Adapun komponen $F$-nya berbeda:
sehingga sebuah proyeksi bergantung pada arah turunnya.}
\end{omfigure}

\section{Keluarga bebas, keluarga pembangun, basis}

\begin{definition}\label{def:b1:vspaces:free}
Sebuah keluarga $(x_1, \dots, x_p)$ berisi vektor $E$ disebut:
\begin{itemize}
  \item \emph{pembangun} (bagi $E$) bila
        $\operatorname{Vect}(x_1,\dots,x_p) = E$;
  \item \emph{bebas}\index{keluarga bebas} (dengan vektornya \emph{bebas
        linear}) bila
        \[
        \lambda_1 x_1 + \dots + \lambda_p x_p = 0
        \implies \lambda_1 = \dots = \lambda_p = 0 ;
        \]
        dan kalau tidak disebut \emph{terikat};
  \item sebuah \emph{basis}\index{basis} bila ia bebas sekaligus pembangun.
\end{itemize}
\end{definition}

\begin{proposition}[Koordinat]\label{prop:b1:vspaces:coordinates}
Keluarga $(e_1, \dots, e_n)$ merupakan \omterm{def:b1:vspaces:free}{basis} $E$ jika dan hanya jika setiap $x \in E$
\emph{secara tunggal} merupakan kombinasi $x = \lambda_1 e_1 + \dots +
\lambda_n e_n$; dan skalar $\lambda_i$-nya menjadi
\emph{koordinat}\index{koordinat} $x$ pada basisnya.
\end{proposition}

\begin{proof}
Pembangun $=$ keberadaan penguraiannya. Ketunggalan $=$
kebebasannya: karena dua penguraian atas $x$ yang sama berselisih sebuah kombinasi
yang sama dengan $0$; lalu kebebasannya memaksa semua koefisiennya --- yaitu
selisih \omterm{prop:b1:vspaces:coordinates}{koordinatnya} --- untuk lenyap. Sebaliknya, sebuah kombinasi
nol yang taksepele memberikan kedua penguraian $0 = \sum \lambda_i
e_i = \sum 0\,e_i$.
\end{proof}

\begin{example}\label{ex:b1:vspaces:bases}
\emph{\omterm{def:b1:vspaces:free}{Basis} kanonik} $K^n$: yaitu $e_i = (0, \dots, 1, \dots, 0)$
(dengan $1$ pada slot $i$). Adapun monomial $(1, X, X^2, \dots, X^n)$: merupakan \omterm{def:b1:vspaces:free}{basis}
$K_n[X]$ (dengan kebebasannya: karena kombinasi nolnya adalah \omterm{def:b1:poly:def}{polinomial} nol, sehingga
semua koefisiennya lenyap, \cref{def:b1:poly:def}). Di dalam
$\C$ atas $\R$: basisnya $(1, \iu)$.
\end{example}

\begin{remark}[Koordinat adalah kerja tim]
\label{rem:b1:vspaces:teameffort}
\omterm{prop:b1:vspaces:coordinates}{Koordinat} pertama $x$ pada \omterm{def:b1:vspaces:free}{basis} $(e_1, \dots, e_n)$
bergantung pada \emph{semua} vektor basisnya, bukan hanya $e_1$. Di dalam
$\R^2$: vektor $(3, 1)$ berkoordinat pertama $3$ pada
\omterm{def:b1:vspaces:free}{basis} kanoniknya, tetapi berkoordinat pertama $2$ pada \omterm{def:b1:vspaces:free}{basis}
$\bigl((1,0), (1,1)\bigr)$ --- pecahkanlah $(3,1) = a(1,0) + b(1,1)$:
maka $b = 1$, $a = 2$. Jadi mengubah satu vektor basisnya mengocok ulang
\emph{setiap} \omterm{prop:b1:vspaces:coordinates}{koordinatnya}; dan \cref{ch:b1:matrices} akan mengemas
pengocokan itu menjadi matriks perubahan \omterm{def:b1:vspaces:free}{basis}.
\end{remark}

\begin{example}[Menguji sebuah calon basis, dari awal sampai akhir]
\label{ex:b1:vspaces:testbasis}
Apakah $\mathcal{F} = (1 + X,\ 1 + X^2,\ X + X^2)$ merupakan \omterm{def:b1:vspaces:free}{basis}
$\R_2[X]$? Tulislah $u_1, u_2, u_3$ bagi ketiga \omterm{def:b1:poly:def}{polinomialnya}.
\emph{Kebebasannya}: sebuah kombinasi nol $a\,u_1 + b\,u_2 + c\,u_3 =
0$ memberikan, koefisien demi koefisien,
\[
a + b = 0, \qquad a + c = 0, \qquad b + c = 0 ;
\]
lalu mengurangkan kedua yang pertama memberikan $b = c$, kemudian yang ketiga memberikan $2b =
0$: sehingga $a = b = c = 0$, jadi \omterm{def:b1:vspaces:free}{bebas}. \emph{Pembangunnya}: alih-alih memecahkan
tiga sistem, perhatikanlah kombinasi setangkupnya
\[
u_1 + u_2 - u_3 = (1 + X) + (1 + X^2) - (X + X^2) = 2 ,
\]
sehingga $1 = \frac12(u_1 + u_2 - u_3)$; lalu
\[
X = u_1 - 1 = \tfrac12\bigl(u_1 - u_2 + u_3\bigr),
\qquad
X^2 = u_2 - 1 = \tfrac12\bigl(-u_1 + u_2 + u_3\bigr).
\]
Monomialnya terletak di \omterm{def:b1:vspaces:span}{rentangnya}, sehingga segalanya demikian: jadi $\mathcal{F}$
merupakan \omterm{def:b1:vspaces:free}{basis}. Sebagai bonusnya, merakit ketiga tampilannya memberikan
\omterm{prop:b1:vspaces:coordinates}{koordinat} sebarang $P = \alpha + \beta X + \gamma X^2$:
\[
P = \frac{\alpha + \beta - \gamma}{2}\,u_1
+ \frac{\alpha - \beta + \gamma}{2}\,u_2
+ \frac{-\alpha + \beta + \gamma}{2}\,u_3 .
\]
(Periksa kewarasannya dengan $P = X$: \omterm{prop:b1:vspaces:coordinates}{koordinatnya} $\bigl(\frac12,
-\frac12, \frac12\bigr)$, sebagaimana ditemukan di atas.) Ada dua pelajarannya: bahwa kesetangkupan
pada keluarganya biasanya menyembunyikan kombinasi jalan pintas; dan bahwa begitu
dimensinya tersedia (\cref{ch:b1:findim}), seluruh paruh
pembangun kerja ini akan datang cuma-cuma --- karena tiga
vektor \omterm{def:b1:vspaces:free}{bebas} pada ruang berdimensi $3$ selalu membentuk \omterm{def:b1:vspaces:free}{basis}.
\end{example}

\begin{proposition}[Kriteria kebebasan yang berguna]\label{prop:b1:vspaces:freecriteria}
\begin{enumerate}
  \item Sebuah keluarga berisi \emph{\omterm{def:b1:poly:def}{polinomial} taknol yang derajatnya berbeda
        berpasangan} bersifat \omterm{def:b1:vspaces:free}{bebas}.
  \item Menambahkan sebuah vektor pada \omterm{def:b1:vspaces:free}{keluarga bebas} menjaganya tetap \omterm{def:b1:vspaces:free}{bebas} jika dan hanya
        jika vektornya berada di luar \omterm{def:b1:vspaces:span}{rentang} keluarganya.
  \item Sebarang subkeluarga sebuah \omterm{def:b1:vspaces:free}{keluarga bebas} bersifat \omterm{def:b1:vspaces:free}{bebas}; dan sebarang keluarga
        yang memuat keluarga pembangun bersifat pembangun.
\end{enumerate}
\end{proposition}

\begin{proof}
(1) Pada sebuah kombinasi nol, lihatlah derajat tertinggi yang hadir: maka
koefisiennya pasti lenyap (karena tak ada yang meniadakan derajat itu), lalu turunlah
secara beruntun.

(2) Jika $x \in \operatorname{Vect}(x_1, \dots, x_p)$, maka relasi $x
- \sum\lambda_i x_i = 0$ bersifat taksepele. Sebaliknya, sebuah kombinasi
nol yang taksepele atas $(x_1, \dots, x_p, x)$ pasti melibatkan $x$ dengan
koefisien yang taknol (kalau tidak ia bertentangan dengan kebebasan keluarga
yang kecil), lalu memecahkannya bagi $x$ menaruhnya di \omterm{def:b1:vspaces:span}{rentangnya}.

(3) Untuk subkeluarganya: sebuah kombinasi nol atas subkeluarganya merupakan kombinasi
seluruh keluarganya dengan koefisien yang hilang ditetapkan $0$; sehingga kebebasan
keluarga besarnya membunuh semuanya. Adapun untuk keluarga besarnya: setiap vektor
$E$ sudah merupakan kombinasi bagian pembangunnya; jadi berikanlah
vektor tambahannya koefisien $0$.
\end{proof}

\begin{example}[Asas tangganya]
\label{ex:b1:vspaces:staircase}
Misalkan $P_0, P_1, \dots, P_n \in K_n[X]$ dengan $\deg P_k = k$ untuk setiap
$k$ (yaitu sebuah ``tangga'' derajat). Maka $(P_0, \dots, P_n)$ merupakan
\omterm{def:b1:vspaces:free}{basis} $K_n[X]$. Adapun kebebasannya adalah
\cref{prop:b1:vspaces:freecriteria}~(1). Sedangkan untuk sifat
pembangunnya, berargumenlah lewat penurunan hingga pada derajatnya: misalkan $Q \in
K_n[X]$, $Q \neq 0$, berderajat $d$, dengan koefisien utamanya $a$,
dan misalkan $b \neq 0$ koefisien utama $P_d$. Maka $Q -
\frac ab P_d$ berderajat $< d$ (karena suku puncaknya saling meniadakan); lalu dengan mengganti
$Q$ dengan selisih ini dan mengulanginya, setelah paling banyak $n + 1$ langkah
kita mencapai \omterm{def:b1:poly:def}{polinomial} nolnya, dan membuka gulungan pengurangannya
menyatakan $Q$ sebagai kombinasi $P_k$-nya. Adapun dua tangga
yang sudah kita temui: pangkat yang tergeser $\bigl((X-a)^k\bigr)_{0 \leq k
\leq n}$ (\cref{exo:b1:vspaces:4}), dan hasil kali Newton
$\bigl((X - x_0)(X - x_1)\cdots(X - x_{k-1})\bigr)_{0 \leq k \leq
n}$, yang dipekerjakan pada soal akhir pekannya.
\end{example}

\begin{example}[Kebebasan pada ruang fungsi]\label{ex:b1:vspaces:functionfree}
Di dalam $\mathcal{F}(\R,\R)$, keluarga $(\eu^{a_1 x}, \dots, \eu^{a_p
x})$ dengan $a_1 < \dots < a_p$ bersifat \omterm{def:b1:vspaces:free}{bebas}: bagilah sebuah kombinasi nol dengan
$\eu^{a_p x}$ lalu biarkan $x \to +\infty$; maka koefisien terakhirnya mati,
dan kita turun secara beruntun (\cref{exo:b1:vspaces:8} memerinci ini beserta
variannya). Jadi kebebasan fungsi dibuktikan lewat \emph{penilaian}:
di titik yang terpilih baik, di tak hingga, atau setelah menurunkannya.
\end{example}

\begin{example}[Sebuah relasi tersembunyi mengecilkan rentang]
\label{ex:b1:vspaces:hiddenrelation}
Di dalam $\mathcal{F}(\R, \R)$, apakah
$\operatorname{Vect}(1,\ \cos^2,\ \sin^2)$? Adapun kesamaan
$\cos^2 + \sin^2 = 1$ merupakan kombinasi nol yang taksepele
\[
1\cdot\mathbf{1} + (-1)\cos^2 + (-1)\sin^2 = 0 :
\]
jadi keluarganya terikat, dan \omterm{def:b1:vspaces:span}{rentangnya} sudah dibangun oleh
$(1, \cos^2)$ saja (karena $\sin^2 = 1 - \cos^2$). Adapun keluarga yang lebih kecil itu
bersifat \omterm{def:b1:vspaces:free}{bebas}: karena $a + b\cos^2 x = 0$ untuk setiap $x$ memberikan, di $x = 0$ dan
$x = \frac\pi2$: $a + b = 0$ dan $a = 0$. Jadi \omterm{def:b1:vspaces:span}{rentangnya} sebuah
\emph{bidang} di dalam ruang fungsinya --- dan ia juga memuat
$\cos 2x = 2\cos^2 x - 1$: sehingga keluarga fungsi trigonometri yang
tampak linear rutin runtuh di bawah kesamaannya,
dan itulah sebabnya kebebasannya harus \emph{dibuktikan}, tak pernah dianggap dari
panjang daftarnya.
\end{example}

\begin{remark}[Jebakan yang lazim]
\label{rem:b1:vspaces:pitfalls}
Ada empat perangkap klasiknya.
\emph{Berpasangan tak mencukupi}: karena di dalam $\R^2$, vektor $(1,0)$,
$(0,1)$, $(1,1)$ tak sebanding berpasangan, namun terikat ---
sebab kebebasan merupakan sifat \emph{seluruh} keluarganya, yang diuji lewat satu
kombinasi yang menyeluruh, tak pernah dua demi dua.
\emph{Vektor nolnya meracuni segalanya}: karena sebarang keluarga yang memuat
$0$ bersifat terikat (sebab $1\cdot 0 = 0$ merupakan relasi yang taksepele), sepolos
apa pun vektor lainnya.
\emph{Gabungan bukan jumlah}: karena $F \cup G$ hampir tak pernah menjadi \omterm{def:b1:vspaces:subspace}{subruang}
(\cref{def:b1:vspaces:subspace}); adapun \omterm{def:b1:vspaces:subspace}{subruang} terkecil
yang memuat keduanya adalah $F + G$, yang biasanya jauh lebih besar daripada gabungannya
--- karena di dalam $\R^2$, dua garis yang berbeda bergabungan menjadi silang, sedangkan berjumlah
seluruh bidangnya.
\emph{Langsung menuntut irisan yang sepele, bukan kelepasan}:
karena dua \omterm{def:b1:vspaces:subspace}{subruang} tak pernah lepas (sebab keduanya memuat $0$); jadi syarat
yang benar adalah $F \cap G = \{0\}$, dan itu harus \emph{dibuktikan},
bukan dibaca dari sebuah gambar --- bandingkanlah
\cref{ex:b1:vspaces:manysupplements}, yang di situ banyak $G$ yang berbeda
berlaku.
\emph{Kebebasan bergantung pada skalarnya}: karena pasangan $(1, \iu)$
\omterm{def:b1:vspaces:free}{bebas} di $\C$ yang dipandang sebagai \omterm{def:b1:vspaces:def}{ruang vektor} atas $\R$, tetapi terikat di $\C$
yang dipandang sebagai \omterm{def:b1:vspaces:def}{ruang vektor} atas $\C$ (sebab $\iu\cdot 1 + (-1)\cdot\iu = 0$).
Jadi ketahuilah selalu \omterm{def:b1:structures:field}{lapangan} mana yang bekerja sebelum menyatakan sebuah \omterm{def:b1:vspaces:free}{keluarga bebas}
--- adapun soal akhir pekan \cref{ch:b1:findim} mengubah persis
kepekaan ini menjadi bukti keirasionalan.
\end{remark}

\begin{remark}[Ke mana bahasa ini pergi]
\label{rem:b1:vspaces:whereused}
Segalanya setelah bab ini berbicara dalam bahasa yang disiapkan di sini.
\cref{ch:b1:findim} mencacah vektor basisnya lalu mengubah ``\omterm{def:b1:vspaces:free}{bebas}'' dan
``pembangun'' menjadi ketaksamaan atas satu bilangan bulat, yaitu
dimensinya. \cref{ch:b1:linmaps} menelaah \omterm{def:b1:logic:map}{pemetaan} yang selaras dengan
kedua operasinya; dan \omterm{def:b1:vspaces:sum}{jumlah langsungnya} menjadi proyektor di sana.
\cref{ch:b1:matrices} menyandikan vektornya lewat \omterm{prop:b1:vspaces:coordinates}{koordinatnya} pada sebuah
\omterm{def:b1:vspaces:free}{basis} --- dengan \cref{prop:b1:vspaces:coordinates} sebagai izin bagi
penyandian itu --- sedangkan \cref{ch:b1:euclid} menambahkan panjang dan sudut
di atas struktur linearnya. Pada jilid Tahun~2 aksioma yang
sama, kata demi kata, berjalan di atas \omterm{def:b1:structures:field}{lapangan} sebarang dan dalam dimensi
tak hingga; karena tak ada apa pun pada bab ini yang memakai kehinggaan di mana pun.
\end{remark}

\begin{remark}[Tiga benang yang diikuti sepanjang Buku 3]
\label{rem:b1:vspaces:threads}
Amatilah tiga gagasan tertentu bab ini bertumbuh.
\emph{Asas tangganya}
(\cref{ex:b1:vspaces:staircase}) muncul kembali sebagai \omterm{def:b1:vspaces:free}{basis} Newton
pada soal akhir pekan bab ini, sebagai \omterm{def:b1:vspaces:free}{basis} binomial
$(B_k)$ di sana, dan sebagai muslihat alternan \omterm{def:b1:poly:def}{polinomial} pada
soal akhir pekan \cref{ch:b1:det}: jadi satu lema, tiga
dividen tanpa determinan.
\emph{Penilaian sebagai uji kebebasan}
(\cref{ex:b1:vspaces:functionfree}) menjadi isomorfisme interpolasi
pada \cref{ch:b1:linmaps}, lalu kriteria Vandermonde
pada \cref{ch:b1:det}, lalu uji Gram pada
\cref{ch:b1:euclid}: jadi refleks yang sama, yang dipertajam tiga kali.
\emph{\omterm{def:b1:vspaces:sum}{Jumlah langsungnya}} (\cref{def:b1:vspaces:sum}) menjadi
proyektor pada \cref{ch:b1:linmaps}, pembelahan ortogonal $E =
F \oplus F^\perp$ pada \cref{ch:b1:euclid}, dan penguraian
terjelaskan-tambah-sisa pada kuadrat terkecil di
soal akhir pekan \cref{ch:b1:multivar}. Jadi sangat sedikit isi buku
ini yang pada dasarnya bukan salah satu dari ketiga gagasan itu dalam pakaian
yang baru.
\end{remark}

\section{Latihan}

\begin{exercise}[$\star$]\label{exo:b1:vspaces:1}
Manakah di antara berikut ini yang merupakan \omterm{def:b1:vspaces:subspace}{subruang}?
\begin{enumerate}
  \item $\{(x, y, z) \in \R^3 : x + 2y - z = 0\}$;
  \item $\{(x, y, z) \in \R^3 : x + 2y - z = 1\}$;
  \item $\{(x, y) \in \R^2 : xy \geq 0\}$;
  \item $\{P \in \R[X] : P(1) = 0\}$;
  \item $\{f \in \mathcal{F}(\R,\R) : f \text{ terbatas}\}$.
\end{enumerate}
\end{exercise}

\begin{exercise}[$\star$]\label{exo:b1:vspaces:2}
Di dalam $\R^3$, apakah $(1, 2, 1)$ termasuk
$\operatorname{Vect}\bigl((1,0,1),\, (1,1,0)\bigr)$? Dan $(2, 1,
1)$? Perikanlah $\operatorname{Vect}\bigl((1,0,1),(1,1,0)\bigr)$ lewat sebuah
persamaan.
\end{exercise}

\begin{exercise}[$\star$]\label{exo:b1:vspaces:3}
Putuskanlah kebebasannya di $\R^3$: $\;\bigl((1,1,0), (1,0,1),
(0,1,1)\bigr)$; $\;\bigl((1,2,3), (2,4,6)\bigr)$;
$\;\bigl((1,0,0), (1,1,0), (1,1,1), (0,1,1)\bigr)$.
\end{exercise}

\begin{exercise}[$\star$]\label{exo:b1:vspaces:4}
Buktikan bahwa $(1, X - 1, (X-1)^2, (X-1)^3)$ merupakan \omterm{def:b1:vspaces:free}{basis} $\R_3[X]$,
lalu berikanlah \omterm{prop:b1:vspaces:coordinates}{koordinat} $X^3$ padanya. \emph{(Lewat Taylor di $1$!)}
\end{exercise}

\begin{exercise}[$\star\star$]\label{exo:b1:vspaces:5}
Di dalam $\R^4$, misalkan $F = \{(x,y,z,t) : x = y = z\}$ dan $G = \{(x,y,z,t)
: x = t = 0\}$. Buktikan bahwa $F \oplus G = \R^4$, lalu uraikanlah
$(1,2,3,4)$ menurutnya.
\end{exercise}

\begin{exercise}[$\star\star$]\label{exo:b1:vspaces:6}
Di dalam ruang barisan, misalkan $F$ \omterm{def:b1:logic:sets}{himpunan} barisan yang
konvergen dan $G = \operatorname{Vect}(u)$ dengan $u_n = (-1)^n$.
Buktikan bahwa $F \cap G = \{0\}$. Apakah $F + G$ seluruh ruang
barisannya?
\end{exercise}

\begin{exercise}[$\star\star$]\label{exo:b1:vspaces:7}
Misalkan $F, G, H$ \omterm{def:b1:vspaces:subspace}{subruang} $E$. Buktikan bahwa
\[
F \cap (G + (F \cap H)) = (F \cap G) + (F \cap H),
\]
lalu tunjukkanlah lewat sebuah contoh di $\R^2$ bahwa kedistributifan
yang tak terbatasi $F \cap (G + H) = (F\cap G) + (F \cap H)$ gagal.
\end{exercise}

\begin{exercise}[$\star\star\star$]\label{exo:b1:vspaces:8}
Buktikan bahwa keluarga $\mathcal{F}(\R, \R)$ berikut bersifat \omterm{def:b1:vspaces:free}{bebas}:
\begin{enumerate}
  \item $(\eu^{a_1 x}, \dots, \eu^{a_p x})$ untuk $a_1 < \dots < a_p$;
  \item $(\cos x, \sin x, \cos 2x, \sin 2x)$;
  \item $(x \mapsto \abs{x - a_1}, \dots, x \mapsto \abs{x - a_p})$
        untuk $a_i$ yang berbeda \emph{(karena sifat \omterm{def:b1:derivative:def}{dapat diturunkannya} gagal tepat di
        satu titik per fungsinya)}.
\end{enumerate}
\end{exercise}

\begin{exercise}[$\star\star\star$]\label{exo:b1:vspaces:9}
Misalkan $E$ \omterm{def:b1:vspaces:def}{ruang vektor} atas $K$ dan $F, G, H$ \omterm{def:b1:vspaces:subspace}{subruang} dengan $F + G =
F + H$, $F \cap G = F \cap H$ dan $G \subseteq H$. Buktikan $G = H$.
Berikanlah contoh penyangkal tanpa hipotesis $G \subseteq H$.
\end{exercise}

\begin{exercise}[$\star\star$]\label{exo:b1:vspaces:10}
Di dalam $\R[X]$, misalkan $\mathcal P$ \omterm{def:b1:logic:sets}{himpunan} \omterm{def:b1:poly:def}{polinomial} yang \emph{genap}
($P(-X) = P(X)$) dan $\mathcal I$ \omterm{def:b1:logic:sets}{himpunan} yang \emph{ganjil}
($P(-X) = -P(X)$). Buktikan bahwa $\R[X] = \mathcal P \oplus \mathcal
I$, lalu tunjukkanlah bahwa $\mathcal P =
\operatorname{Vect}(1, X^2, X^4, \dots)$, yakni bahwa \omterm{def:b1:poly:def}{polinomial}
yang genap tepat merupakan \omterm{def:b1:poly:def}{polinomial} dalam $X^2$.
\end{exercise}

\begin{exercise}[$\star\star$]\label{exo:b1:vspaces:11}
Misalkan $(x_1, x_2, x_3)$ \omterm{def:b1:vspaces:free}{keluarga bebas} sebuah \omterm{def:b1:vspaces:def}{ruang vektor} real $E$.
Buktikan bahwa $(x_1 + x_2,\ x_2 + x_3,\ x_3 + x_1)$ bersifat \omterm{def:b1:vspaces:free}{bebas}. Apakah
keluarga serupa yang berisi empat vektor $(x_1 + x_2,\ x_2 + x_3,\ x_3 +
x_4,\ x_4 + x_1)$ \omterm{def:b1:vspaces:free}{bebas} ketika $(x_1, x_2, x_3, x_4)$ demikian?
\end{exercise}

\begin{exercise}[$\star\star\star$]\label{exo:b1:vspaces:12}
Misalkan $E$ \omterm{def:b1:vspaces:def}{ruang vektor} atas $\R$ (atau $\C$) dan $F_1, \dots,
F_k$ \omterm{def:b1:vspaces:subspace}{subruang} \emph{sejati} $E$ (dengan setiap $F_i \neq E$).
\begin{enumerate}
  \item Tanganilah kasus $k = 2$ secara langsung: jika $F_1 \not\subseteq
        F_2$ dan $F_2 \not\subseteq F_1$, pilihlah $x \in F_1
        \setminus F_2$ dan $y \in F_2 \setminus F_1$ lalu carilah letak $x
        + y$.
  \item Buktikan secara umum bahwa $E \neq F_1 \cup \dots \cup F_k$: yakni bahwa
        \omterm{def:b1:vspaces:def}{ruang vektor} atas \omterm{def:b1:structures:field}{lapangan} yang tak hingga tak pernah menjadi gabungan
        hingga \omterm{def:b1:vspaces:subspace}{subruang} sejati. \emph{(Ambillah $k$ yang minimal, pilihlah
        $x \in F_1$ di luar $F_i$ lainnya, pilihlah $y \notin F_1$,
        lalu ikutilah garis $t \mapsto y + tx$.)}
\end{enumerate}
\end{exercise}

\section{Soal: interpolasi, tiga basis bagi satu ruang}

\begin{problem}\label{pb:b1:vspaces:1}
Tetapkanlah $n + 1$ titik yang \emph{berbeda} $x_0, x_1, \dots, x_n$ pada $\R$.
Soal ini menengok kembali \omterm{thm:b1:poly:lagrange}{interpolasi Lagrange}
(\cref{thm:b1:poly:lagrange}) dengan mata bab ini: bahwa
ruang $\R_n[X]$ membawa tiga \omterm{def:b1:vspaces:free}{basis} yang alami --- yaitu milik Lagrange,
milik Newton, dan (untuk titik yang berjarak sama) \omterm{def:b1:vspaces:free}{basis} binomialnya ---
dan setiap basisnya membuat satu pertanyaan menjadi mudah. Jalannya berakhir pada sebuah teorema
aritmetika yang sungguhan: yaitu pencirian Pólya atas \omterm{def:b1:poly:def}{polinomial}
yang memetakan $\Z$ ke dalam $\Z$.\index{interpolasi Lagrange}
\index{selisih terbagi}\index{polinomial bernilai bulat}

\medskip
\noindent\textbf{Bagian I --- \omterm{def:b1:vspaces:free}{Basis} Lagrangenya.} Untuk $0 \leq i
\leq n$ tetapkanlah
\[
L_i \;=\; \prod_{j \neq i} \frac{X - x_j}{x_i - x_j} \;\in\;
\R_n[X].
\]
\begin{enumerate}
  \item Periksalah bahwa $\deg L_i = n$ dan bahwa $L_i(x_j) = 1$ bila $j =
        i$, dan $0$ bila $j \neq i$.
  \item Buktikan bahwa keluarga $(L_0, \dots, L_n)$ bersifat \omterm{def:b1:vspaces:free}{bebas}.
  \item Buktikan bahwa untuk setiap $P \in \R_n[X]$,
        \[
        P \;=\; \sum_{i=0}^{n} P(x_i)\, L_i ,
        \]
        lalu simpulkanlah bahwa $(L_0, \dots, L_n)$ merupakan \omterm{def:b1:vspaces:free}{basis}
        $\R_n[X]$. \emph{(Tinjaulah selisih kedua ruasnya
        lalu cacahlah akarnya, \cref{cor:b1:poly:nroots}.)}
  \item Simpulkanlah teorema interpolasinya: bahwa untuk sebarang nilai $y_0,
        \dots, y_n \in \R$ ada $P \in
        \R_n[X]$ yang \emph{tunggal} dengan $P(x_i) = y_i$ untuk setiap $i$. Pada \omterm{def:b1:vspaces:free}{basis}
        Lagrangenya, apakah \omterm{prop:b1:vspaces:coordinates}{koordinat} sebuah \omterm{def:b1:poly:def}{polinomial} $P$?
  \item Buktikanlah kesamaan
        \[
        \sum_{i=0}^{n} L_i = 1
        \qquad\text{dan, untuk } 0 \leq k \leq n,\qquad
        \sum_{i=0}^{n} x_i^{k}\, L_i = X^{k} .
        \]
\end{enumerate}

\medskip
\noindent\textbf{Bagian II --- \omterm{def:b1:vspaces:free}{Basis} Newton dan \omterm{pb:b1:vspaces:1}{selisih
terbaginya}.} Tetapkanlah $N_0 = 1$ dan $N_k = (X - x_0)(X - x_1) \cdots
(X - x_{k-1})$ untuk $1 \leq k \leq n$. Untuk sebuah fungsi $f$ yang terdefinisi
di simpulnya, definisikanlah \emph{\omterm{pb:b1:vspaces:1}{selisih terbaginya}} lewat
$f[x_i] = f(x_i)$ dan
\[
f[x_i, \dots, x_{i+k}] \;=\;
\frac{f[x_{i+1}, \dots, x_{i+k}] - f[x_i, \dots, x_{i+k-1}]}
{x_{i+k} - x_i} .
\]
\begin{enumerate}[resume]
  \item Buktikan bahwa $(N_0, N_1, \dots, N_n)$ merupakan \omterm{def:b1:vspaces:free}{basis}
        $\R_n[X]$.
  \item Hitunglah $f[x_0, x_1]$ dan $f[x_0, x_1, x_2]$ dalam
        nilai $f$, lalu hitunglah semua \omterm{pb:b1:vspaces:1}{selisih terbagi}
        $f(x) = x^2$ pada tiga simpul yang sebarang.
  \item (Lema Aitken) Misalkan $R$ menginterpolasi $f$ di $x_0, \dots,
        x_{n-1}$ dan $Q$ menginterpolasi $f$ di $x_1, \dots, x_n$,
        keduanya berderajat $\leq n - 1$. Buktikan bahwa
        \[
        S \;=\; \frac{(X - x_0)\,Q - (X - x_n)\,R}{x_n - x_0}
        \]
        menginterpolasi $f$ di $x_0, x_1, \dots, x_n$.
  \item Simpulkanlah, lewat induksi pada cacah simpulnya, bahwa
        koefisien $X^{k}$ pada interpolan $f$ di $x_0,
        \dots, x_k$ tepat sama dengan $f[x_0, \dots, x_k]$.
  \item Buktikanlah \emph{rumus interpolasi Newton}: bahwa
        interpolan $f$ di $x_0, \dots, x_n$ adalah
        \[
        P \;=\; \sum_{k=0}^{n} f[x_0, \dots, x_k]\, N_k ,
        \]
        lalu turunkanlah rumus \omterm{def:b1:topology:closed}{tertutupnya}
        \[
        f[x_0, \dots, x_k] \;=\; \sum_{i=0}^{k}
        \frac{f(x_i)}{\prod_{j \neq i,\, j \leq k} (x_i - x_j)} ,
        \]
        yang menunjukkan bahwa $f[x_0, \dots, x_k]$ tak bergantung pada
        urutan simpulnya.
\end{enumerate}

\medskip
\noindent\textbf{Bagian III --- Simpul yang berjarak sama: operator
selisihnya.} Mulai sekarang simpulnya adalah $0, 1, 2, \dots$ dan, untuk sebuah
\omterm{def:b1:poly:def}{polinomial} $P$, kita menetapkan
\[
\Delta P(X) = P(X + 1) - P(X),
\qquad
B_k = \frac{X(X-1)\cdots(X-k+1)}{k!} \quad (B_0 = 1).
\]
\begin{enumerate}[resume]
  \item Tunjukkan bahwa jika $\deg P = m \geq 1$ dengan koefisien utamanya
        $a$, maka $\deg \Delta P = m - 1$ dengan koefisien utamanya
        $m\,a$, dan bahwa $\Delta$ membunuh konstantanya.
  \item Tunjukkan bahwa $(B_0, B_1, \dots, B_n)$ merupakan \omterm{def:b1:vspaces:free}{basis}
        $\R_n[X]$ dan bahwa $\Delta B_k = B_{k-1}$ untuk $k \geq 1$.
  \item (Rumus selisih maju Newton) Buktikan bahwa setiap $P
        \in \R_n[X]$ memenuhi
        \[
        P \;=\; \sum_{k=0}^{n} \bigl(\Delta^{k} P\bigr)(0)\, B_k .
        \]
  \item Buktikan bahwa untuk setiap $k \geq 0$,
        \[
        \bigl(\Delta^{k} P\bigr)(0) \;=\;
        \sum_{j=0}^{k} (-1)^{k-j} \binom{k}{j} P(j) .
        \]
  \item Tunjukkan bahwa jika $\deg P = n$ dengan koefisien utamanya $a_n$,
        maka $\Delta^{n} P$ merupakan konstanta $n!\,a_n$ dan
        $\Delta^{n+1} P = 0$.
\end{enumerate}

\medskip
\noindent\textbf{Bagian IV --- \omterm{pb:b1:vspaces:1}{Polinomial bernilai bulat}.} Sebuah
\omterm{def:b1:poly:def}{polinomial} $P \in \R[X]$ disebut \emph{bernilai bulat} bila $P(m) \in
\Z$ untuk setiap $m \in \Z$.
\begin{enumerate}[resume]
  \item Buktikan bahwa setiap $B_k$ bernilai bulat. \emph{(Tanganilah $m
        \geq k$, $0 \leq m < k$ dan $m < 0$ secara terpisah; adapun untuk $m =
        -q < 0$, tunjukkanlah $B_k(-q) = (-1)^k \binom{q + k - 1}{k}$.)}
  \item Buktikanlah \emph{pencirian Pólya}: bahwa $P \in \R_n[X]$
        bernilai bulat jika dan hanya jika \omterm{prop:b1:vspaces:coordinates}{koordinatnya} pada
        \omterm{def:b1:vspaces:free}{basis} $(B_0, \dots, B_n)$ berupa bilangan bulat.
  \item Simpulkanlah: bahwa jika $P \in \R_n[X]$ mengambil nilai bulat di $n + 1$
        bilangan bulat yang \emph{berurutan} $a, a+1, \dots, a+n$, maka $P$
        bernilai bulat. \emph{(Geserlah: terapkanlah telaahnya pada $Q(X)
        = P(X + a)$.)}
  \item Simpulkanlah dari pertanyaan 16 bahwa hasil kali $k$ bilangan bulat
        yang berurutan selalu terbagi oleh $k!$.
  \item Misalkan $P = \dfrac{X(X+1)(2X+1)}{6}$. Hitunglah tabel Newtonnya
        di $0, 1, 2, 3$, tulislah $P$ pada \omterm{def:b1:vspaces:free}{basis} $(B_k)$, lalu
        simpulkanlah bahwa $P$ bernilai bulat meskipun tak ada satu pun dari
        koefisien monomialnya bulat. Periksalah $\Delta P =
        (X+1)^2$ lalu simpulkanlah $P(m) = 1^2 + 2^2 + \dots + m^2$ untuk
        $m \in \N$.
\end{enumerate}

\medskip
\noindent\textbf{Bagian V --- Dividennya.}
\begin{enumerate}[resume]
  \item Misalkan $P \in \R_n[X]$ menginterpolasi nilai $2^i$ di $i =
        0, 1, \dots, n$. Tunjukkan bahwa $P = B_0 + B_1 + \dots + B_n$
        dan bahwa $P(n + 1) = 2^{n+1} - 1$: sehingga ``pola
        pelipatduaannya'' selalu patah tepat pada titik berikutnya.
  \item (Antiturunan yang diskret) Buktikan bahwa untuk setiap bilangan bulat $m
        \geq 1$ dan $k \geq 0$,
        \[
        \sum_{j=0}^{m-1} B_k(j) \;=\; B_{k+1}(m),
        \]
        yakni kesamaan tongkat hoki $\sum_{j=k}^{m-1}
        \binom{j}{k} = \binom{m}{k+1}$.
  \item Jabarkanlah $X^2$ dan $X^3$ pada \omterm{def:b1:vspaces:free}{basis} $(B_k)$ lalu simpulkanlah
        rumus \omterm{def:b1:topology:closed}{tertutup} bagi $\sum_{j=0}^{m-1} j^2$ dan
        $\sum_{j=0}^{m-1} j^3$; lalu pulihkanlah kesamaan Nicomachus
        $1^3 + \dots + m^3 = (1 + \dots + m)^2$.
  \item Ambillah $n = 2$ dengan simpul $0, 1, 2$. Tulislah \omterm{prop:b1:vspaces:coordinates}{koordinat}
        $X^2$ pada ketiga \omterm{def:b1:vspaces:free}{basis} soal ini: yaitu \omterm{def:b1:vspaces:free}{basis}
        monomialnya, \omterm{def:b1:vspaces:free}{basis} Lagrangenya, dan \omterm{def:b1:vspaces:free}{basis} Newtonnya. Periksalah ketiga
        jawabannya terhadap pertanyaan 4 dan 9.
  \item Sintesis. Dalam empat kalimat: gagasan \omterm{def:b1:vspaces:def}{ruang vektor} mana
        yang membuat pertanyaan 4 otomatis; mengapakah \omterm{def:b1:vspaces:free}{basis} Newton menghitung
        \omterm{prop:b1:vspaces:coordinates}{koordinatnya} \emph{secara rekursif} sedangkan \omterm{def:b1:vspaces:free}{basis} Lagrange
        membacanya \emph{seketika}; kriteria kebebasan mana
        yang dibagi kedua basisnya; dan dalam arti persis apa teorema
        Pólya mengatakan bahwa kebulatan sebuah \omterm{def:b1:poly:def}{polinomial} merupakan
        sifat \omterm{prop:b1:vspaces:coordinates}{koordinatnya} \emph{pada \omterm{def:b1:vspaces:free}{basis} yang tepat}.

\end{enumerate}
\end{problem}
